% Part: lambda-calculus
% Chapter: lambda-definability
% Section: lists

\documentclass[../../../include/open-logic-section]{subfiles}

\begin{document}

\olfileid{lam}{ldf}{lst}
\olsection{فہرستاں}

فہرست $[M_1, M_2, \ldots, M_n]$ دی تعریف ایویں کیتی جا سکدی اے:
\[
  [M_1, M_2, \ldots, M_n] \ident \lambd[sz][s M_1 (s M_2 (\ldots (s M_n z)))]
\]
عام بولی وچ، فہرست دی نمائندگی کرن والا ٹرم اوہدا «قدر جوڑدے جان والا» فنکشن اے۔ ایہ دو ٹرم لیندا اے۔ پہلا فنکشن اے جیہڑا ہُن تک جوڑی قدر تے اک نواں رکن لے کے نویں جوڑی قدر موڑدا اے۔ دوجا جوڑ شروع کرن والی قدر اے۔ فہرست والا فنکشن اک اک کر کے رکن جوڑدا اے تے آخر حاصل موڑدا اے۔

\begin{digress}
  جے تسی فنکشنی پروگرامنگ توں واقف او، تاں ایہ تعریف \emph{سجّے پاسے دی تہہ} ورگی لگے گی۔ فہرست نوں اوہدے رکناں دی سجّے پاسے دی تہہ والے فنکشن وجوں تعریف کیتا جاندا اے۔
\end{digress}

ایس کوڈ بندی دی بنیاد اُتے اسیں کجھ کم آؤن والے فنکشن سوکھے ای تعریف کر سکدے آں:
\begin{align*}
  \fn{Sum} & \ident
  \lambd[l][l \, (\lambd[xa][\fn{Add} \, x \, a])\,  \num{0}]\\
  \fn{Len} & \ident
  \lambd[l][l \, (\lambd[xa][\fn{Add} \, a \, \num{1}]) \num{0}]
\end{align*}

$\fn{Sum}$ چرچ دیاں عددی عبارتاں دی فہرست دا جوڑ کڈھدا اے۔ ایہ فہرست اُتے اک اک رکن جوڑ کے کم کردا اے۔ شروعاتی قدر $\num{0}$ اے تے ہر رکن لئی نویں قدر وجوں $\fn{Add} \, x\, a$ موڑدا اے (جتھے $x$ ہُن والا رکن اے تے $a$ ہُن تک جوڑی قدر اے)۔ حاصل سارے رکناں دا جوڑ اے۔

\begin{prob}
  $\fn{Len}$ کیہ کردا اے؟ سمجھاؤ۔
\end{prob}

\begin{prob}
  اک فنکشن تعریف کرو جیہڑا فہرست دی ترتیب اُلٹا دے۔
\end{prob}
\end{document}
